\begin{houseviewbox}[最终结论]
文章最有价值的部分，是提醒投资者不要把所有科技硬件放进同一篮子，也不要把创新药 BD 和电力建设简单视为情绪对手盘。它最薄弱的部分，是把未经复核的资金数字和跨层级的成本传导，推成了确定的交易结论。

在证据允许的范围内，半导体设备是值得优先补证的产业方向，创新药 BD 是需要逐资产拆分的现金流方向，电力设备是需要用订单到现金桥检验的趋势方向。三者当前均不足以从主题直接跃迁为个股投资结论。研究行动为：维持六个样本的卫星观察；只在下一份财报、订单/验收、临床/会计或招投标/现金证据闭环时重新开启估值；在此之前不做无差别追涨。
\end{houseviewbox}

\section{研究状态}
本 PDF 是完整的研究叙事和证据审计交付，但不是通过 R1--R4 全部估值与投资委员会门禁的可投资报告。原因是核心池为空而非遗漏：这是在证据不足时应有的研究结论。
